\documentclass[12pt,a4paper]{article}

\usepackage[english]{babel}
\usepackage[utf8]{inputenc}
\usepackage{amsmath,amsfonts}
\usepackage{graphicx}
\usepackage[margin=1in]{geometry}
\usepackage{enumitem}
\usepackage{comment}
\usepackage{setspace}
\usepackage{hyperref}

\title{\textbf{TITLE OF THE MANUSCRIPT}}

\author{\textbf{TITLE IN A DIFFERENT LANGUAGE}}

\date{}

\begin{document}

\maketitle


\textbf{Type of Contribution }

(Theoretical Essay, Pilot Study, Project Registration, Experience Report, Research Report, Literature Review, Tutorial, according to the \href{https://cadernos.abralin.org/index.php/cadernos/about/submissions}{Journal Guidelines}  ).
\vspace{0.5cm}
%this for manual spacing; feel free to change the number per your needs

\textbf{Author’s Full Name}\footnote{The first author is the lead author. The author’s footnote must have: degree, affiliation, e-mail, ORCID ID, and the roles of each author according to \href{http://www.abralin.org/circulares/cadlin/CRediT-CadLin.pdf}{CRediT} taxonomy in case of co-authorship.}

\textbf{Author’s Full Name} \footnote{From the second author on, all of the other contributors are co-authors. The co-authors’ footnote must have: degree, affiliation, e-mail, ORCID ID, and the role of each author according to  \href{http://www.abralin.org/circulares/cadlin/CRediT-CadLin.pdf}{CRediT}  taxonomy.}

\vspace{0.5cm}
%this for manual spacing; feel free to change the number per your needs

\textbf{ABSTRACT}: The text of the ABSTRACT must be typed in the same font and size as the text, written in the language of the manuscript, with only the first letter capitalized, without bold, in only one paragraph of text, containing 150 to 250 words. When appropriate, it should present the objective(s) of the research, the methodology, and the main results. 
\vspace{0.1cm}
%this for manual spacing; feel free to change the number per your needs

\textbf{KEYWORDS}: There must be 3 to 5 keywords separated by a dot, written in the language of the manuscript, and typed in the same font and size as the text, with only the first letter capi-talized.
\vspace{0.5cm}
%this for manual spacing; feel free to change the number per your needs

\textbf{ABSTRACT}: ABSTRACT in a different language of the manuscript, following the same characteristics described above for the abstract in the language of the manuscript. The Bra-zilian Association of Linguistics is a signatory of the Helsinki Initiative on Multilingualism in Scientific Communication. The signatories of this initiative support the promotion of lin-guistic diversity through the publication of research results in all languages. 
\vspace{0.1cm}
%this for manual spacing; feel free to change the number per your needs

\textbf{KEYWORDS}: KEYWORDS in a different language from the manuscript, following the same characteristics described above for the keywords in the language of the manuscript.

\setstretch{1.5}

\section{Main text} 
The manuscript must be typed in Word (or a compatible program), must have be-tween 12 and 30 pages, including references, and must be formatted as follows: A4 page layout, no line/page number, 2.5cm margins; the text must be in one column, in 12pt Times New Roman font, with a 1.5cm space between lines (except for the References, which should have a 1cm space between lines) and with a 1.25cm first line indentation (except in the Ref-erences and in the headings of the sections, which must be left-justified and have no indenta-tion).

Italics (\textit{itálico}) should be used only to indicate words or phrases in a language differ-ent from the manuscript. “Quotation marks” should be used only for short quotations (up to three lines) and to highlight words or phrases, example sentences, titles, slang, or neolo-gisms. 

Notes must be written as footnotes in 10pt font and should only be restricted to in-formation relevant to the text. They should not contain references.\footnote{The notes should be placed as footnotes rather than endnotes and must be restricted to content, not to citations. Times New Roman font, size 10.}


CadLin receives submissions in any language. In the case of texts written in lan-guages other than Portuguese and English, authors should indicate a list of at least five pos-sible referees. Before submitting the manuscripts to CadLin, the authors should make a care-ful textual revision. The editors and/or reviewers of the Journal will not make any textual revision.

Authors who submit their paper to CadLin will automatically declare that ethical principles have been observed, as well as the legal requirements related to ethics in research, both in the country where they maintain their institutional bond and in the country where the people participating in the research live.

Papers accepted for publication in CadLin may be accompanied by badges from the Center for Open Science that recognize open scientific practice: (a) publicly available data, (b) publicly available materials, and (c) pre-registered research plans. If you would like more information on how to obtain an open science practice badge, please refer to the fol-lowing: http://www.abralin.org/circulares/cadlin/badges-cadlin.pdf. To request the inclusion of one of these badges, authors should complete the Application for Open Scientific Practic-es Badges, available here: http://www.abralin.org/circulares/cadlin/badges-form-cadlin.pdf.

\section{About the sections} 
Section headings must be numbered (except for the Acknowledgments, if any, and the References). They must be in bold with just the first letters of the heading and any prop-er nouns capitalized. They should be followed by a blank line.

\section{Illustrations (Tables, Graphs, Figures, Boxes)} 
The illustrations (Tables, Graphs, Figures, Tables) will have their legends and con-tents typed in the text font, size 10, without bold. The legend (title) of each illustration should begin with the name of the illustration to which it refers (in bold), numbered sequen-tially (for example: \textbf{Table 1., Table 2., Graph 1., Graph 2.}), and inserted below the illustra-tion. The source should be cited (in case of previously published illustrations). All illustra-tions should be inserted in the body of the text immediately after the paragraph in which they were cited, and not in the form of annexes. Please try to use grayscale illustrations whenever possible. If you have to use colored illustrations, please make sure that they are \href{https://jfly.uni-koeln.de/color/index.html}{friendly to colorblind people.} 

\section{Citations} 
Quotations must follow the ABNT NBR 10520 guidelines. In indirect citations (para-phrases), the author should be cited in parentheses by last name, in all capital letters, sepa-rated by a comma from the year of publication (SWERTS, 1997). If the name of the author is cited in the text, only the year should be indicated in parentheses: According to Oliveira Jr (2000), […].

Whenever the indication of the page is necessary, as in a direct quotation, the number of the page must be placed right after the year, separated from it by a comma and preceded by “p.”, for example, (OLIVEIRA JR, 2000, p. 95). Direct quotations of up to three lines are made inside the text in double quotation marks, while direct quotations of more than three lines should be separated from the text with a 4cm indentation to the left, font size 10, with-out quotation marks and without indentation of the first line.

The citations of several works by the same author published in the same year should be distinguished by lower case letters after the date, without spacing (FERREIRA, 2007a).

When the work has two or three authors, all may be indicated, separated by semico-lons (HUETTIG; ROMMERS; MEYER, 2011); when there are more than three authors, the first name is indicated, followed by \textbf{et al.} (ALMEIDA \textbf{et al.}, 2013). In this case, it is advisa-ble to indicate all authors in the References.

The References at the end of the text must also comply with ABNT norms. The con-tent of the articles and the accuracy of the references are the sole responsibility of the au-thors. The Editors and the Brazilian Association of Linguistics do not assume any responsi-bility for the opinions or statements of the authors. All and only the works of authors cited in the text should appear in the References, which are typed with single spacing between lines, separated by a simple space, and listed in alphabetical order by the surname of the first author. Where available, the authors should include the URLs and DOIs of the refer-ences used. 

\section*{Acknowledgements} 

Acknowledgments and mention of research funding (scholarships, financial aid, and the respective funding agencies) that enabled the research should not be included in a foot-note but in an Acknowledgments section at the end of the article, before the References.

\vspace{2cm}


\section*{References} 

\setlength{\parindent}{0em}
%this command ensures that there is no tab spacing in the References section

\textbf{References: (according to ABNT NBR 6023 (2018), the DOI must be referenced when-ever possible). 
}

BRAYNER, A. R. A.; MEDEIROS, C. B. Incorporação do tempo em SGBD orientado a objetos. \textit{In}: SIMPÓSIO BRASILEIRO DE BANCO DE DADOS, 9., 1994, São Paulo. Anais [...]. São Paulo: USP, 1994. p. 16-29. (\textbf{Example of a work published in Proceedings})
\vspace{0.5cm}


DANTAS, José Alves \textit{et al}. Regulação da auditoria em sistemas bancários: análise do cenário internacional e fatores determinantes. \textbf{Revista Contabilidade & Finanças}, São Paulo, v. 25, n. 64, p. 7-18, jan./abr. 2014. DOI http://dx.doi.org/10.1590/S1519-70772014000100002. Acesso em: 20 maio 2014 \textbf{(Example of an article published in a journal with a DOI}).
\vspace{0.5cm}

LUCK, Heloisa. \textbf{Liderança em gestão escolar.} 4. ed. Petrópolis: Vozes, 2010. (\textbf{Example of a book})
\vspace{0.5cm}

MORAES, João Antônio de; RILLIARD, Albert. “Prosody and Emotion in Brazilian Portu-guese”. \textit{In}: ARMSTRONG, Meghan E.; HENRIKSEN, Nicholas; VANRELL, Maria del Mar. \textbf{Intonational Grammar in Ibero-Romance}: Approaches across linguistic subfields. Amsterdam: John Benjamins, 2016, p. 135-152. (Example of a book chapter)

RODRIGUES, Ana Lúcia Aquilas. \textbf{Impacto de um programa de exercícios no local de trabalho sobre o nível de atividade física e o estágio de prontidão para a mudança de comportamento}. 2009. Dissertação (Mestrado em Fisiopatologia Experimental) – Faculdade de Medicina, Universidade de São Paulo, São Paulo, 2009. (\textbf{Example of an aca-demic work – thesis, dissertation, etc.})

SEKEFF, Gisela. O emprego dos sonhos. \textbf{Domingo}, Rio de Janeiro, ano 26, n. 1344, p. 30-36, 3 fev. 2002. (\textbf{Example of a journal article without URL or DOI})



\end{document}